%% kapa: exam=1 round=33 year=2022 period=2 form=A subject=appraisal_law subject_ko=감정평가관계법규 seq=19 number=19
%% source: Q-Net 감정평가사 기출문제 (cst00309) artlSeq=5209997 file=raw/1008/5209997/2022년도 제33회 감정평가사 1차시험 2교시(감정평가 관계법규, 회계학).pdf
%% conversion: kordoc 4.24.51 -> tools/gen1.py (auto); review=proofread-sonnet; flags=none
\begin{question}{19}
부동산 가격공시에 관한 법령상 표준지공시지가에 관한 설명으로 옳은 것을 모두 고른 것은?
\begin{framedtext}
ㄱ. 표준지공시지가의 공시에는 표준지 및 주변토지의 이용상황이 포함되어야 한다. \par
ㄴ. 표준지공시지가는 일반적인 토지거래의 지표가 된다. \par
ㄷ. 도시개발사업에서 환지를 위하여 지가를 산정할 때에는 표준지공시지가를 기준으로 하지 아니한다. \par
ㄹ. 최근 1년간 시·군·구별 지가변동률이 전국 평균 지가변동률 이하인 지역의 표준지에 대해서는 하나의 감정평가법인등에 의뢰하여 표준지공시지가를 조사·평가할 수 있다. \par
\end{framedtext}
\begin{choices}
  \item ㄱ, ㄴ
  \item ㄱ, ㄷ
  \item ㄱ, ㄴ, ㄹ
  \item ㄴ, ㄷ, ㄹ
  \item ㄱ, ㄴ, ㄷ, ㄹ
\end{choices}
\end{question}
\answer{1}
